\section{Introduction}

The companion write-up in this repository proved that the stationary equilibrium of
\citet{Aiyagari1994} is unique at his logarithmic calibration, and located exactly what stands in
the way for relative risk aversion $\mu\in\{3,5\}$. Every obstacle it met traces to the fixed point
of the infinite-horizon problem: a condition on the response of consumption to the interest rate
fails in a tail of wealth that the stationary distribution reaches; an induction that would prove
the response bound has a step multiplier $1/\Thorn>1$, so it cannot close as a supremum argument;
and the only floor on capital supply comes from an ergodic identity that loses a factor
$\|u'\|^2/\operatorname{Var}$. Overlapping generations remove all three at once. A household
that lives $J$ periods solves a finite backward induction from an explicit terminal condition, so
a multiplier $1/\Thorn$ per step compounds to $\Thorn^{-J}$, which at Aiyagari's numbers is about
$1.03$. Wealth is bounded by $J$ periods of accumulation, so any condition needs checking on a
known bounded region. And the distribution of the youngest cohort is common to every interest
rate, so an ordering of policies pushes forward to an ordering of the whole age profile of
distributions by induction on age, with no Doeblin atom and no uniqueness of a stationary
distribution.

The target of the programme is uniqueness of the stationary equilibrium of a stochastic life-cycle
economy in the manner of \citet{Huggett1996} at $\mu\in\{3,5\}$. It is reached at $\mu=1$ and not
above, and the paper is as much a record of where it is blocked as of what is proved.

Sections~\ref{sec:diamond} to~\ref{sec:eqm} are the foundations. Section~\ref{sec:diamond} does
Diamond's two-period economy, where the steady state is unique for every risk aversion and the
household-side comparative static that blocks the general case runs the wrong way without doing any
harm. Section~\ref{sec:lc} sets out the life-cycle household, the finite-horizon sandwich, and the
propensity floor. Section~\ref{sec:ret} adds retirement, a pension and a means test, and settles
which policy a means test makes non-monotone and in which state space. Section~\ref{sec:eqm} proves
existence of a stationary equilibrium with both boundary conditions discharged from primitives rather
than assumed: capital supply in a life cycle is a finite composition of finite sums, so existence
needs no stationary distribution and no compactness argument beyond Berge, and it can be bracketed by
two affine bounds on a cohort's accumulated assets sharing a slope, with every coefficient an explicit
recursion in the primitives. At Aiyagari's calibration the hypotheses hold on an interval of about
five percentage points containing the true rate of $5.09\%$.

Sections~\ref{sec:thm1} to~\ref{sec:inc} are about uniqueness, and they are best read as three
attempts on one obstacle. Uniqueness needs capital supply to rise with the interest rate, and the
only route the literature offers is Light's Theorem 1, that saving rises with the rate at every
state. Section~\ref{sec:thm1} reports how far that goes: it holds at every age below a critical
relative risk aversion of about $3.4$, so the life cycle covers Aiyagari's $\mu=3$ and not his
$\mu=5$; the Euler reduction of it is false at logarithmic utility exactly where an equilibrium sits,
and a different sufficient condition, on the value gap between two rates, is true there instead
(Section~\ref{sec:log}). The section also separates the mechanism from the usual story: the reversal
happens in the certainty-equivalent household, which faces no risk at all, so precautionary saving is
not what drives it, and we give the criterion in closed form and prove its certainty-equivalent half
for every risk aversion.

Section~\ref{sec:mfg} asks what mean field games offer. The Lasry--Lions argument was dismissed for
the stationary Bewley model because it needs a common initial distribution, which overlapping
generations have; so the monotonicity condition becomes a question to check rather than one to
dismiss. Date by date it fails, at every risk aversion, and not for the expected reason. What the
argument actually uses is weaker --- only the summed pairing is signed --- and that weaker condition
holds throughout the bracket in which the equilibrium is known to lie. But pushed to its weakest form
it collapses to monotone capital supply, so the route is not logically weaker than the one it was
meant to bypass. The section is a closure, with the reduction and the two theorems it produces kept
because they are about the life-cycle household rather than about mean field games.

Section~\ref{sec:inc} is the route that works at $\mu=1$. Single crossing does not ask capital supply
to rise, only to fall more slowly than the firm's demand; over the proved bracket the demand curve is
steep enough that a proof may lose more than the true increment and still close. Coupling the two
economies on a common earnings path and splitting the step at the lower-rate state gives an envelope
in which the failure of Theorem 1 is priced rather than excluded, and its per-stage multiplier is
damped by the propensity floor. Two further facts make the estimate small enough: a cohort's earnings
marginal is the invariant law at every age, so the mass weighting is exact in the coordinate where the
failure lives; and the multiplier is one plus the variance of the risk tilt, an identity rather than a
bound. That variance is certified below what the damping allows, in exact rational arithmetic, at every
stage and over the whole asset range. What that delivers is Theorem 1 at every state and age, hence
monotone capital supply, hence uniqueness --- at logarithmic utility, at one calibration, with the last
hypothesis discharged by certificate rather than by proof term. Section~\ref{sec:prog} records what
remains.

As in the companion, every result stated with a formal name has been checked by the Lean kernel
against Mathlib; the development is at
\url{https://github.com/LeanEconomics/LeanEconomics}, directory \texttt{LeanEconomics/OLG}, and
the numerical work is in \texttt{WriteUps/WriteUpOLG/numerics}.
